\chapter{Discussion}

\section{Strengths of the Current Prototype}

The strongest aspect of the current DaQuVa prototype is that the language now
covers a complete cleaning workflow rather than isolated syntax examples. A
user can load CSV or SQLite data, run scans, combine filters with
\texttt{AND}/\texttt{OR}, apply string predicates, inspect the most relevant
rows with \texttt{sort} and \texttt{limit}, detect near duplicates, merge them,
and export the cleaned table. The \texttt{summarize} command rounds out this
workflow with an aggregate proof table that reports before/after row counts for
each stage, so a reviewer can confirm that cleaning actually reduced the data
rather than just inspecting individual rows. The new REPL strengthens the
workflow further because it supports incremental exploration without repeatedly
running a full script.

The implementation also keeps destructive operations visible. Database writes
require explicit syntax such as \texttt{allowDanger}, and the SQLite output path
uses a safer temporary-table replacement strategy. This matches the original
safety-first design goal of the language.

\section{Limitations and Future Work}

The prototype is still intentionally small. The expression language supports
comparisons, boolean composition, and a few string predicates, but it does not
yet provide joins, date operations, or user-defined validators. Aggregation is
currently limited to the fixed-schema \texttt{summarize} proof table; there is
no general \texttt{group by} with user-chosen aggregate functions. The REPL is useful for exploration, but it does not yet provide
autocomplete, command history persistence, or static diagnostics before
execution.

Future work should focus on broadening the type system, adding a formal test
suite around the demo scenarios, and extending SQL pushdown to more operations.
Another important direction is improving documentation synchronization so that
the BNF, parser implementation, examples, and report remain aligned as the DSL
evolves.
